%%%%%%%%%%%%%%%%%%%%%%%%%%%%%%%%%%%%%%%%%%%%%%%%%%%%%%%%%%%%%%%%%
                % Probabilistic transformer %
%%%%%%%%%%%%%%%%%%%%%%%%%%%%%%%%%%%%%%%%%%%%%%%%%%%%%%%%%%%%%%%%%


As mentioned in Chapter 1, there is an unexplored area in the study of transformers expressiveness; a study of non-determinism\bigcifu{non-determinism} is missing.

When $\TF$ is a deterministic transformer, $\TF^i(\alpha)$ is an exact symbol for all $\alpha \in \Sigma^*$ and $i \in \mathbb{N}$.
On the contrary, when the token selection process of the transformer samples the distribution, it becomes a random variable.
Therefore, we introduce the following definition.


\begin{definition}[Probabilistic transformer]
    We call probabilistic transformer to the transformers that use the probabilistic selector.
    We say that a family of probabilistic transformers $\{\TF_n\}_{n \in \mathbb{N}}$ recognizes a language $\gL$ when: 
    \[
    \alpha \in \gL \iff \Pr[\TF_{|\alpha|}^*(\alpha) = \qaccept] > 2/3
    \]
\end{definition}

\bigcifu{Estaría bueno agregar unas palabras comentando qué cosas no estamos modelando. Por ejemplo, la normalización de los vectores al final de cada capa interna o cosas asi. Por lo menos para que se note que sabemos que este modelo no es perfecto.}